% GENERATED FILE. Edit canonical lessons, then run npm run generate:books.
% canonical-source-hash: fnv1a64:846fb0d85c1e0df0
% canonical-lessons: TE-C209-durada, TE-C209-narsu, TE-C209-pariksa, TE-C209-vaidyam, TE-C209-nuduru

\chapter{Itch, Nurse, Test, Medical treatment, Forehead}
\label{ch:te-itch-nurse-test-medical-treatment-forehead}
\hlchaptermodality{209}

\begin{chapteropening}
\textbf{By the end of this chapter:} \emph{I can say an itch, a nurse, a test, medical treatment and the forehead.}

Complete the last lesson of chapter 209: I can say an itch, a nurse, a test, medical treatment and the forehead.
\end{chapteropening}

\section[\texorpdfstring{\te{దురద} (durada)}{దురద (durada)}]{\te{దురద} (durada) — an itch}

\label{lesson:TE-C209-durada}



\begin{quote}
Before the new one: say the Telugu for to be born, then the Telugu for to confront.
\end{quote}

\subsection*{You'll want to know: \te{దురద}}
\textbf{\te{దురద}} — \emph{durada} — \textquotedblleft{}an itch\textquotedblright{}.

\begin{grammarlens}[title={What you've built}]
One more word for this chapter.
\end{grammarlens}

\subsection*{Your turn}
\begin{itemize}
\raggedright
  \item \emph{Say it:} \emph{durada}
  \item \emph{Say it:} \emph{durada}, once more
  \item \emph{Say it:} say \emph{durada}
  \item \emph{Recall:} say the Telugu for to be born, then the Telugu for to confront, then say \emph{durada} again
\end{itemize}

\begin{tcolorbox}[breakable,colback=teal!4,colframe=teal!35!black,title={Before you move on}]
What does \te{దురద} mean? (\textquotedblleft{}An itch\textquotedblright{}.) Say it once more.
\end{tcolorbox}
\section[\texorpdfstring{\te{నర్సు} (narsu)}{నర్సు (narsu)}]{\te{నర్సు} (narsu) — a nurse}

\label{lesson:TE-C209-narsu}



\begin{quote}
Before the new one: say the Telugu for to confront, then the Telugu for an itch.
\end{quote}

\subsection*{You'll want to know: \te{నర్సు}}
\textbf{\te{నర్సు}} — \emph{narsu} — \textquotedblleft{}a nurse\textquotedblright{}.

\begin{grammarlens}[title={What you've built}]
One more word for this chapter.
\end{grammarlens}

\subsection*{Your turn}
\begin{itemize}
\raggedright
  \item \emph{Say it:} \emph{narsu}
  \item \emph{Say it:} \emph{narsu}, once more
  \item \emph{Say it:} say \emph{narsu}
  \item \emph{Recall:} say the Telugu for to confront, then the Telugu for an itch, then say \emph{narsu} again
\end{itemize}

\begin{tcolorbox}[breakable,colback=teal!4,colframe=teal!35!black,title={Before you move on}]
What does \te{నర్సు} mean? (\textquotedblleft{}A nurse\textquotedblright{}.) Say it once more.
\end{tcolorbox}
\section[\texorpdfstring{\te{పరీక్ష} (parīkṣa)}{పరీక్ష (parīkṣa)}]{\te{పరీక్ష} (parīkṣa) — a test, an examination}

\label{lesson:TE-C209-pariksa}



\begin{quote}
Before the new one: say the Telugu for an itch, then the Telugu for a nurse.
\end{quote}

\subsection*{You'll want to know: \te{పరీక్ష}}
\textbf{\te{పరీక్ష}} — \emph{parīkṣa} — \textquotedblleft{}a test, an examination\textquotedblright{}.

\begin{grammarlens}[title={What you've built}]
One more word for this chapter.
\end{grammarlens}

\subsection*{Your turn}
\begin{itemize}
\raggedright
  \item \emph{Say it:} \emph{parīkṣa}
  \item \emph{Say it:} \emph{parīkṣa}, once more
  \item \emph{Say it:} say \emph{parīkṣa}
  \item \emph{Recall:} say the Telugu for an itch, then the Telugu for a nurse, then say \emph{parīkṣa} again
\end{itemize}

\begin{tcolorbox}[breakable,colback=teal!4,colframe=teal!35!black,title={Before you move on}]
What does \te{పరీక్ష} mean? (\textquotedblleft{}A test, an examination\textquotedblright{}.) Say it once more.
\end{tcolorbox}
\section[\texorpdfstring{\te{వైద్యం} (vaidyaṁ)}{వైద్యం (vaidyaṁ)}]{\te{వైద్యం} (vaidyaṁ) — medical treatment}

\label{lesson:TE-C209-vaidyam}



\begin{quote}
Before the new one: say the Telugu for a nurse, then the Telugu for a test.
\end{quote}

\subsection*{You'll want to know: \te{వైద్యం}}
\textbf{\te{వైద్యం}} — \emph{vaidyaṁ} — \textquotedblleft{}medical treatment\textquotedblright{}.

\begin{grammarlens}[title={What you've built}]
One more word for this chapter.
\end{grammarlens}

\subsection*{Your turn}
\begin{itemize}
\raggedright
  \item \emph{Say it:} \emph{vaidyaṁ}
  \item \emph{Say it:} \emph{vaidyaṁ}, once more
  \item \emph{Say it:} say \emph{vaidyaṁ}
  \item \emph{Recall:} say the Telugu for a nurse, then the Telugu for a test, then say \emph{vaidyaṁ} again
\end{itemize}

\begin{tcolorbox}[breakable,colback=teal!4,colframe=teal!35!black,title={Before you move on}]
What does \te{వైద్యం} mean? (\textquotedblleft{}Medical treatment\textquotedblright{}.) Say it once more.
\end{tcolorbox}
\section[\texorpdfstring{\te{నుదురు} (nuduru)}{నుదురు (nuduru)}]{\te{నుదురు} (nuduru) — the forehead}

\label{lesson:TE-C209-nuduru}



\begin{quote}
Before the new one: say the Telugu for a test, then the Telugu for medical treatment.
\end{quote}

\subsection*{You'll want to know: \te{నుదురు}}
\textbf{\te{నుదురు}} — \emph{nuduru} — \textquotedblleft{}the forehead\textquotedblright{}.

\begin{grammarlens}[title={What you've built}]
One more word for this chapter.
\end{grammarlens}

\subsection*{Your turn}
\begin{itemize}
\raggedright
  \item \emph{Say it:} \emph{nuduru}
  \item \emph{Say it:} \emph{nuduru}, once more
  \item \emph{Say it:} say \emph{nuduru}
  \item \emph{Recall:} say the Telugu for a test, then the Telugu for medical treatment, then say \emph{nuduru} again
\end{itemize}

\begin{tcolorbox}[breakable,colback=teal!4,colframe=teal!35!black,title={Before you move on}]
What does \te{నుదురు} mean? (\textquotedblleft{}The forehead\textquotedblright{}.) Say it once more.
\end{tcolorbox}
